% 2012 年北京卷文科第 4 题：输出 S 的程序框图。
% 来源 PDF：普通高考/2012/2012北京文.pdf，第 1 页（印面 719）左栏。
% 流程：开始 → k=0,S=1 → 判断 k<3 → 是：S=S·2^k，k=k+1，回到判断；否：输出 S → 结束。
\begin{tikzpicture}[semithick, >=Stealth,
  block/.style={draw, minimum width=3.0cm, minimum height=0.62cm, align=center},
  dec/.style={draw, diamond, aspect=2.6, inner sep=1pt, align=center},
  term/.style={draw, rounded corners=6pt, minimum width=1.5cm, minimum height=0.52cm},
  io/.style={draw, trapezium, trapezium left angle=70, trapezium right angle=110,
             minimum width=2.2cm, minimum height=0.58cm},
  arr/.style={->}]
  \node (start) [term] at (0,6.4) {开始};
  \node (init)  [block] at (0,5.4) {$k=0,\ S=1$};
  \node (dec)   [dec]  at (0,3.6) {$k<3$?};
  \node (s)     [block] at (3.6,3.6) {$S=S\cdot 2^{k}$};
  \node (k)     [block] at (3.6,4.4) {$k=k+1$};
  \node (out)   [io]   at (0,2.3) {输出 $S$};
  \node (end)   [term] at (0,1.3) {结束};
  \draw[arr] (start) -- (init);
  \draw[arr] (init) -- (dec);
  \draw[arr] (dec) -- node[above] {是} (s);
  \draw[arr] (s) -- (k);
  \draw[arr] (k.west) -- (0,4.4);
  \draw[arr] (dec) -- node[right] {否} (out);
  \draw[arr] (out) -- (end);
\end{tikzpicture}
